\section{Release: Datasets, Benchmark, Access, and Maintenance}
\label{sec:release}

\paragraph{What is released.}
(i)~The generator, under the MIT license, with the full descriptor libraries,
knob registry, calibration pipeline, validation suite, command-line interface,
and a browser front end for interactive scenario construction. (ii)~Generated
datasets under CC~BY~4.0, described by a datasheet
\citep{gebru2021datasheets} and machine-readable Croissant metadata
\citep{akhtar2024croissant}. (iii)~Sixty benchmark scenario configurations
spanning nine sites, six building types, thirteen driver cohorts (three
fitted from real corpora, one labeled fixture, nine hand-specified), and six
equipment mixes, including sweep families for fleet scale, infrastructure,
flexibility share, climate/season, and tariff structure. (iv)~A reference
corpus: ten heterogeneous buildings on one campus $\times$ twelve months
$\times$ 150 samples per building-month --- 18{,}000 dataset units, each an
independent weather realization (per-sample dry-bulb, solar, dew-point, and
wind draws re-simulated through EnergyPlus) with the clean noise profile.

\paragraph{Corpus integrity, stated fully.}
Every unit passes all hard validation invariants (zero errors across 18{,}000
units). Soft advisories are disclosed rather than suppressed --- 11{,}920 of
the 18{,}000 units carry at least one --- typically warnings for charger-concurrency near saturation (a scenario-design property)
and for DR programs whose season excludes the simulated month; the released
corpus documentation includes the complete warning profile. One recording
subtlety is worth stating because it misled our own internal audit: weather
realization is configured \emph{per building}, so the batch-level manifest
field reports no global weather profile; the authoritative record of each
unit's weather draw is the per-unit configuration file, which pins the exact
offsets applied.

\paragraph{Benchmark scope.}
The released scheduling baseline evaluates unidirectional (V1G) smart
charging. We scope it so deliberately: dispatch of bidirectional discharge
and stationary storage is a \emph{decision} of the controller under
evaluation, exactly as the dataset ships storage specifications rather than
dispatch schedules; publishing a reference V2B dispatch would entangle the
benchmark with one controller's choices. The dataset itself carries every
channel a V2B controller requires (bidirectional charger rates, storage
specifications, DR commitments, tariffs), and the shift pairs of
\S\ref{sec:utility} are released as learning benchmarks alongside the
scheduling baseline.

\paragraph{Access and maintenance.}
The generator is developed in a public repository; the reference corpus is
deposited with a persistent identifier (DOI) \owed{Zenodo DOI, WS-F}.
Dataset versioning is the triplet (generator revision, configuration, seed):
any released unit regenerates bitwise from its manifest, so the corpus is
recoverable from the generator alone. We commit to maintaining the repository
and the deposit through the review period and beyond; issues and corrections
are tracked publicly.

\paragraph{Compute.}
One dataset unit (one building-month) costs $49$\,s on average
(${\sim}117$\,s when the EnergyPlus load cache is cold), dominated by the
building simulation; the 18{,}000-unit corpus (19.5\,GB) took
${\approx}246$ CPU-hours, ${\approx}8.2$\,h wall-clock on 30 workers.
Single-scenario generation for interactive use completes in roughly two
minutes.
